\VerdierPhysicalPageAnchor{0079}%
\VerdierSemanticAnchor{ch1-2-5-7-morphisms}
\noindent\hspace*{2em}--\ Soient $X$ et $Y$ deux objets de $K(C)$,
$f,g\colon X\rightrightarrows Y$ deux morphismes de
$\operatorname{comp}(C)$. La relation ``$f$ et $g$ sont homotopes'' est une
relation d'équivalence sur
$\operatorname{Hom}_{\operatorname{comp}(C)}(X,Y)$, et on définit
$\operatorname{Hom}_{K(C)}(X,Y)$ comme le quotient de
$\operatorname{Hom}_{\operatorname{comp}(C)}(X,Y)$ par cette relation
d'équivalence.

\VerdierSemanticAnchor{ch1-2-5-7-composition}
\noindent\hspace*{2em}--\ Il résulte immédiatement de (2.5.5) que la loi de
composition sur les morphismes de $\operatorname{comp}(C)$ passe au quotient
et définit une loi de composition sur les morphismes de $K(C)$.

\VerdierSemanticAnchor{ch1-2-5-7-additive-category}
Il résulte immédiatement de (2.5.5) qu'on a bien défini ainsi une catégorie et
que cette catégorie est additive.

\VerdierSemanticAnchor{ch1-2-5-7-quotient-functor}
Le foncteur canonique
$\operatorname{comp}(C)\xrightarrow{Q}K(C)$ qui est l'identité sur les objets
et qui associe à tout morphisme $f$ de $\operatorname{comp}(C)$ sa classe
d'équivalence $\dot f$ est additif.

\VerdierSemanticAnchor{ch1-2-5-7-homotopy-class}
Le morphisme $\dot f$ de $K(C)$ image du morphisme $f$ de
$\operatorname{comp}(C)$ par le foncteur canonique est appelé la
\emph{classe de $f$ modulo homotopie}.

\VerdierSemanticAnchor{ch1-2-5-7-name}
La catégorie $K(C)$ est appelée la \emph{catégorie des complexes de $C$ à
homotopie près}.

\VerdierSemanticAnchor{ch1-2-5-8}
\noindent\textbf{2.5.8.}\quad
Le foncteur canonique $\operatorname{comp}(C)\xrightarrow{Q}K(C)$ possède la
propriété suivante : tout foncteur de $K(C)$ dans une catégorie $A$ fournit,
par composition avec $Q$, un foncteur de $\operatorname{comp}(C)$ dans $A$,
d'où, pour toute catégorie $A$, un foncteur :

\VerdierSemanticAnchor{ch1-2-5-8-hom-functor}
\begin{equation*}
  \mathcal H\!om(Q,A)\colon
  \mathcal H\!om(K(C),A)
  \longrightarrow
  \mathcal H\!om(\operatorname{comp}(C),A)
\end{equation*}
($\mathcal H\!om$ désigne la catégorie des foncteurs).

\VerdierSemanticAnchor{ch1-2-5-8-fully-faithful}
Ce foncteur est pleinement fidèle, injectif sur les objets et l'image est la
sous-catégorie pleine de :

\VerdierSemanticAnchor{ch1-2-5-8-essential-image}
\begin{equation*}
  \mathcal H\!om(\operatorname{comp}(C),A)
\end{equation*}
définie par les foncteurs de $\operatorname{comp}(C)$ dans $A$ qui
transforment tout couple de morphismes homotopes de $\operatorname{comp}(C)$
en morphismes égaux.

\VerdierSemanticAnchor{ch1-2-5-8-uniqueness}
La catégorie $K(C)$ munie du foncteur canonique
$\operatorname{comp}(C)\xrightarrow{Q}K(C)$ est déterminée à isomorphisme
unique près par cette propriété.

\VerdierSemanticAnchor{ch1-2-5-9}
\noindent\textbf{2.5.9.}\quad
Il résulte immédiatement de (2.5.5) et de (2.5.8) qu'il existe un et un seul
foncteur de translation sur $K(C)$ tel que le foncteur canonique :

\VerdierSemanticAnchor{ch1-2-5-9-translation}
\begin{equation*}
  \operatorname{comp}(C)\longrightarrow K(C)
\end{equation*}
commute (strictement) aux translations.

\VerdierSemanticAnchor{ch1-2-5-9-strict-z-category}
Ce foncteur de translation est un automorphisme de $K(C)$ et par suite définit
sur $K(C)$ une structure de $\mathbb Z$-catégorie stricte.

\VerdierSemanticAnchor{ch1-2-5-9-strict-z-functor}
Le foncteur canonique $\operatorname{comp}(C)\longrightarrow K(C)$ est un
$\mathbb Z$-foncteur strict.
